% =====================================================================
%  4.3 Derivasjonsregler for polynom — ekstra utfordring
% =====================================================================

\begin{infoboks}{Velg metode selv}
Arket bygger på 4.3. Vis regning og begrunnelse. Skriv tangenter på formen $y = mx + c$.
Gang ut produkter før du deriverer. Husk enhet når oppgaven har en kontekst.
\end{infoboks}

\nivaahode{2}

\oppg{U1}{2}{\mh}
Kostnaden ved å lage $x$ enheter er $K(x) = 0,01x^2 + 20x + 5000$ kroner (tenkt modell).
\begin{tasks}(1)
  \task Finn $K'(x)$ og regn ut $K'(100)$. Forklar hva tallet betyr, med enhet.
  \task Regn ut $K(101) - K(100)$. Hvorfor er svaret nesten lik $K'(100)$?
\end{tasks}

\oppg{U2}{2}{\uh}
Finn tangentene til $f(x) = x^3 - 6x$ som er parallelle med linja $y = 6x - 1$.

\oppg{U3}{2}{\uh}
La $f(x) = x^2 - 4x + 6$. Tangenten i $x = 1$ danner sammen med de to aksene en trekant.
Finn arealet av trekanten.

\oppg{U4}{2}{\prog}
Programmet regner ut stigningstallet til tangenten med en symmetrisk differansekvotient.
\begin{kode}
def f(x):
    return x**4 - x

def derivert(f, a, h=0.0001):
    return (f(a + h) - f(a - h)) / (2 * h)

a = 1
m = derivert(f, a)
c = f(a) - m * a
print("m =", round(m, 2), "og c =", round(c, 2))
\end{kode}
\begin{tasks}(1)
  \task Kjør programmet. Hva skriver det ut? Hva betyr $m$ og $c$?
  \task Finn tangenten i $x = 1$ eksakt med derivasjonsreglene. Stemmer svaret med programmet?
  \task Endre $f$ til $x^3 - 2x$ og $a$ til $2$. Finn tangenten med programmet og kontroller med reglene.
\end{tasks}

\nivaahode{3}

\oppg{U5}{3}{\uh}
Ola vil finne $a$ og $b$ slik at $f(x) = ax^2 + bx$ har $f'(1) = 4$ og $f(2) = 10$. Han finner ingen løsning.
\begin{tasks}(1)
  \task Forklar hvorfor det ikke finnes noen løsning.
  \task For hvilken verdi av $f(2)$ ville det fungert? Hvor mange løsninger er det da?
\end{tasks}

\oppg{U6}{3}{\uh}
Vis at tangentene til $y = x^2$ i $x = a$ og $x = -a$ ($a \neq 0$) skjærer hverandre på $y$-aksen.
Finn skjæringspunktet.

\oppg{U7}{3}{\uh}
Punktet $(1, -3)$ ligger utenfor grafen til $y = x^2$. Finn alle tangenter til grafen som går gjennom punktet.
\hint{Tangenten må røre grafen et sted. Kall det stedet $a$.}

\begin{oppsummering}
Gang ut og forkort før du deriverer. Tangenten i $x = a$ bruker tallet $f'(a)$ og punktet $(a, f(a))$.
Når et krav gir ett stigningstall, løs $f'(x) = k$.
\end{oppsummering}
